\section{投资问题：低PE是真便宜，还是峰值错觉}

恒逸2026E看似只有约7倍PE，但分母取法、H2正常化、少数股东和现金转化决定“便宜”是否成立。核心判断是，周期调整PE比单一DCF更适合当前阶段，基本面概率锚为15.6398元；再加入市场锚14.86元和归一化Street锚16.3130元，得到模型值15.6575元，发布目标15.7元。相对15.06元盘中价上行约4.3\%，安全边际不足以支持积极评级。

主方法不是把一份券商79.1亿元预测乘更高倍数，而是使用独立House 77.5亿元、38.80694189亿股已披露稀释分母和18.746155亿元员工计划现金。转债有效利息与员工计划股份支付未量化，主模型不回加。

\section{三情景基本面价值}

熊／基准／牛情景全年归母为65／77.5／95亿元，对应H2 7.5／20.0／37.5亿元；PE为6／7.5／9倍。情景价值均先计算总权益价值，加上员工计划最高现金，再除以已披露稀释股本。概率为25\%／55\%／20\%。

\begin{exhibitbox}[三情景权益价值]
\noindent\begin{tabularx}{\textwidth}{L{1.8cm}R{1.8cm}R{1.8cm}R{1.8cm}R{1.4cm}R{1.7cm}R{1.3cm}X}
\toprule
\textbf{情景} & \textbf{全年归母} & \textbf{H2归母} & \textbf{稀释EPS} & \textbf{PE} & \textbf{价值} & \textbf{概率} & \textbf{核心含义} \\
\midrule
熊 & 65.0 & 7.5 & 1.675 & 6.0$\times$ & 10.53元 & 25\% & 价差回落、现金与资本折价 \\
基准 & 77.5 & 20.0 & 1.997 & 7.5$\times$ & 15.46元 & 55\% & H1真实、H2正常化 \\
牛 & 95.0 & 37.5 & 2.448 & 9.0$\times$ & 22.52元 & 20\% & 高价差、国内和现金共振 \\
\bottomrule
\end{tabularx}
\end{exhibitbox}
\sourcenote{估值审计，模型截止2026-07-22。价值均为$[\text{\songti 归母}\times PE+18.746155]\div38.80694189$；金额单位为亿元，股本为亿股。}

概率加权基本面锚为：
\[
15.6398=10.5328\times25\%+15.4611\times55\%+22.5152\times20\%.
\]
核心公允区间按House基准利润的6.5--8.0倍权益价值计算，为13.46--16.46元，发布简化为13.5--16.5元。10.5--22.5元是尾部情景，不是日常目标区间；将22.5元写成“合理目标”会把低概率牛情景伪装成基准。

\begin{keyinsight}[三情景的“所以呢”]
当前价15.06元接近基准情景15.46元，交易的不是低谷利润，而是多数2026反转。\textbf{研究上，13.5元以下且基本面未证伪时进入条件升级区；16.5元以上只有在H2利润和现金流同时上修时才维持正向研究结论。}
\end{keyinsight}

\section{市场锚：约束模型脱离成交现实}

截至7月21日五个完整交易日，成交量加权均价为14.6555元；7月22日13:52盘中价为15.06元。市场锚取两者中点14.8578元，发布四舍五入为14.86元。该锚不是技术分析给基本面背书，而是在预告驱动的高波动阶段防止模型用假精确远离真实成交。

市场锚权重15\%，低于基本面与Street。若价格短期上冲但财务证据不变，市场锚会小幅提高目标，却不足以机械转为买入；若价格下跌且基本面未失效，市场锚也会限制目标的下移幅度。真正改变评级的仍是利润、现金和资本结构。

\section{Street锚：先归一股本，再给予有限权重}

国信原始PDF给出2026--2028归母79.1／78.0／85.6亿元，目标区间16.5--17.6元，中点17.05元，估值方法为FCFE与可比PE。其预测表采用约36.03亿股，低于主模型38.80694189亿股。直接把17.05元作为目标会忽略员工计划和股本差异。

本报告先把国信中点转换为总权益价值，再加入同一员工计划现金，最后除以主模型股本：
\[
16.3130=\frac{17.05\times36.03+18.746155}{38.80694189}.
\]
得到归一化Street锚16.3130元，权重20\%。其余七份公开摘要缺少原始PDF和完整方法，权重为零，不能与国信组成伪“均值”。

\begin{exhibitbox}[多锚目标价桥]
\noindent\begin{tabularx}{\textwidth}{L{3.35cm}R{2.0cm}R{1.7cm}X}
\toprule
\textbf{锚} & \textbf{每股价值} & \textbf{权重} & \textbf{作用} \\
\midrule
House基本面概率锚 & 15.6398元 & 65\% & 独立盈利、周期情景和稀释股本 \\
市场隐含锚 & 14.86元 & 15\% & 五日VWAP与盘中价中点 \\
归一化Street锚 & 16.3130元 & 20\% & 一份原始PDF，经股本／现金归一 \\
模型值 & 15.6575元 & 100\% & 三项加权 \\
发布目标 & 15.7元 & -- & 控制假精确 \\
\bottomrule
\end{tabularx}
\end{exhibitbox}
\sourcenote{HY-MKT-KLINE；HY-MKT-QUOTE；HY-BRK-GUOXIN；analysis/valuation\_audit.md，截止2026-07-22。}

最终模型为：
\[
15.6575=15.6398\times65\%+14.86\times15\%+16.3130\times20\%.
\]
模型值相对15.06元上行3.97\%；发布目标15.7元相对当前价上行4.25\%，报告表述约4.3\%。详细小数用于复算，读者行动使用15.7元。

\section{当前价隐含利润：市场已经付了多少}

在7.5倍PE下，按当前经济股本34.3972亿股，15.06元隐含归母69.07亿元，为House基准89.1\%；加入转债潜在股份、不回加有效利息，隐含74.90亿元，为96.6\%；再加入员工计划上限并扣除员工现金，隐含75.42亿元，为97.3\%。后两项是估值敏感性，不是会计if-converted EPS。

\begin{exhibitbox}[当前价格隐含的2026归母]
\noindent\begin{tabularx}{\textwidth}{L{4.2cm}R{2.1cm}R{2.0cm}X}
\toprule
\textbf{状态} & \textbf{隐含归母} & \textbf{占House} & \textbf{解读} \\
\midrule
经济股本 & 69.07亿元 & 89.1\% & 看似仍有利润缺口 \\
转债股数敏感性 & 74.90亿元 & 96.6\% & 稀释后接近基准 \\
已披露稀释并扣员工现金 & 75.42亿元 & 97.3\% & 主口径下已计入大部分基准 \\
\bottomrule
\end{tabularx}
\end{exhibitbox}
\sourcenote{HY-MKT-QUOTE；HY-OFF-CB；HY-OFF-ESOP7；analysis/valuation\_audit.md，盘中价截至2026-07-22 13:52。}

这组反推是中性评级的关键。投资者若在15.06元买入，主要不是等待77.5亿元兑现，而是押注95亿元牛情景、9倍PE或现金质量显著上修。赔率与已知利润增速并不对称。

\section{PB／ROE与FCFE交叉检查}

2026Q1归母权益约271.19亿元。按经济股本计算每股净资产约7.88元；按已披露稀释股本并加入员工计划现金约7.47元。当前价对应约1.9--2.0倍PB。若77.5亿元归母兑现，短期ROE很高，但它包含峰值周期利润，不能用成长股PB线性外推。

国信FCFE中枢约17.79元，高于本报告15.7元，主要差异来自持续价差、资本开支、股本稀释、股份支付和Street口径归一。DCF／FCFE对远期裂解、二期投入和终值高度敏感，本报告只把其作为上边界检查，不作为唯一锚。

PEG／PSG权重为零。2026同比增速来自2025低基数和周期价差，不是稳定单位销量、ASP和结构性市场份额扩张；用PEG会把周期反转误写成成长。

\begin{riskbox}[估值失效条件]
若H2归母低于18亿元、现金流未修复或资本开支／融资上升，利润与PE应同时下调；若H2达到35亿元以上且现金覆盖资本开支，牛情景可上调。\textbf{行动上，目标价随证据更新，不因股价触及目标而反向修改假设。}
\end{riskbox}
